\scanpage{47}
\unsure{satisfont} aux conditions précédentes $\unsure{\alpha)}$, donc une structure à puissances divisées canonique sur $pA$.

\uline{Exemple} 3 \quad Soit $A$ intègre, \deleted{\unsure{à valuation}} tel que les seules \deleted{\unsure{son corps des fractions soit}} \unsure{car.} résiduelles soient $p$ et $0$. Soit $\pi$ un élément tel que les caractéristiques résiduelles de $A/\pi A$ soient $p$. À quelle condition peut-on trouver des $\unsure{\pi^{(n)}|\pi^{[n]}}$ ?

Supposons que $A$ soit de valuation discrète. Il faut écrire \deleted{$\unsure{\forall}$} $\pi^n/n!\in A$, i.e.
\[
v(\pi^n/n!)\geq 0\qquad\text{pour }\forall n
\]
i.e.
\[
n v(\pi)-v(n!)\geq0\qquad\text{pour tout }n.
\]
Soit \deleted{$\unsure{v_p}$} $v(p)=e$ l'indice de ramification absolue de $A$, \unsure{on a} \deleted{$v(N)$} pour $\forall$ entier $N$,
\[
v(N)=e\,v_p(N).
\]
\deleted{Soit de plus $\rho=v(\pi)$.} La condition s'écrit
\[
\deleted{\unsure{\rho}}v(\pi)\geq\deleted{\unsure{e}}e\frac{\deleted{\unsure{e}}v_p(n!)}{n}\deleted{\unsure{e}}\qquad\text{pour tout }n
\]
i.e.
\[
v(\pi)\geq e\lambda(p)\qquad\text{où}\quad\lambda(p)=\sup_n\frac{v_p(n!)}{n}\inserted{=\frac1{p-1}}
\]
\margin{\unsure{Legendre}}
est une constante ne dépendant que de $p$, et qui est $\leq1$ \ldots\ Donc on trouve la condition
\margin{\deleted{\unsure{déterminer} $\lambda(p)$}}
\[
v(\pi)\geq\frac{e}{\unsure{r|p}-1}.
\]
Supposons que $\pi$ soit l'uniformisante, donc $v(\pi)=1$, on trouve la condition
\[
\boxed{e\leq p-1}
\]
\margin{46}
